\documentclass[11pt]{article}
\usepackage[margin=0.82in]{geometry}
\usepackage{amsmath,amssymb,mathtools}
\usepackage{booktabs,tabularx,array}
\usepackage[dvipsnames]{xcolor}
\usepackage[colorlinks=true,linkcolor=NavyBlue,urlcolor=BrickRed,citecolor=NavyBlue]{hyperref}
\usepackage{microtype}
\usepackage{fancyhdr}
\usepackage{enumitem}

\definecolor{ink}{HTML}{17212B}
\definecolor{accent}{HTML}{315B7D}
\definecolor{pale}{HTML}{EDF4F8}
\definecolor{passgreen}{HTML}{16794A}
\color{ink}
\pagestyle{fancy}
\fancyhf{}
\lhead{Proctor odd symplectic group}
\rhead{Proof and verification}
\cfoot{\thepage}
\setlength{\headheight}{14pt}
\setlength{\parindent}{0pt}
\setlength{\parskip}{5pt}
\newcommand{\Sfun}{\mathcal S}
\newcommand{\Sym}{\operatorname{Sym}}
\newcommand{\Sp}{\operatorname{Sp}}
\newcommand{\rad}{\operatorname{rad}}
\newcommand{\Spec}{\operatorname{Spec}}
\newcommand{\PASS}{\textcolor{passgreen}{\textbf{PASS}}}
\newcommand{\summarybox}[1]{\begin{center}\fcolorbox{accent}{pale}{\parbox{0.91\linewidth}{#1}}\end{center}}

\title{\textbf{Proctor's Odd Symplectic Group}\\[3pt]
\large What can be averaged, exact replacements, and computational verification}
\author{Proof and reproducibility note}
\date{17 July 2026}

\begin{document}
\maketitle

\summarybox{\textbf{Outcome.} The proposed normalized Haar expectation for the
full complex odd symplectic group is not defined: the group contains both a
vector group and $\mathbb C^\times$. Three valid replacements were proved and
tested. Exact constant-term calculations verify the two compact formulas
through total degree six. The full algebraic invariant series is exactly $1$.
For finite fields, four matrix groups of orders $24$, $432$, $12{,}000$, and
$11{,}520$ were enumerated in dimensions three and five; all 38 tested
coefficients agree between direct permutation-matrix averages and the proposed
cycle formula. Selected coefficients also agree with explicit orbit counts.}

\section{The obstruction comes before computation}

Let $V=W\oplus L$ over $\mathbb C$, where $\dim W=2n$, $L=\mathbb C e_0$,
and $\omega$ is a nondegenerate alternating form on $W$. Extend it by
\[
 \Omega((w,s),(w',s'))=\omega(w,w'),\qquad \rad\Omega=L.
\]
Consider the full stabilizer
\[
 G=\{g\in GL(V):\Omega(gv,gw)=\Omega(v,w)\}.
\]
Because the radical is intrinsic, $gL=L$. In the decomposition $W\oplus L$,
every element therefore has the form
\begin{equation}
 g=\begin{pmatrix}A&0\\ \ell&a\end{pmatrix},\qquad
 A\in\Sp(W),\quad \ell\in W^*,\quad a\in\mathbb C^\times.
\end{equation}
Conversely, every matrix in (1) preserves $\Omega$. Thus
\begin{equation}
 G\cong W^*\rtimes\bigl(\Sp_{2n}(\mathbb C)\times\mathbb C^\times\bigr).
\end{equation}
The vector group $W^*$ is the unipotent radical, so $G$ is nonreductive; both
$W^*$ and $\mathbb C^\times$ are noncompact. A locally compact group has a
normalized Haar probability exactly when it is compact. Consequently
\[
 \frac{1}{\operatorname{vol}(G)}\int_G f(g)\,dg
\]
is not a defined normalization. This is a mathematical obstruction, not a
numerical convergence problem. A finite box cutoff would depend on coordinates
and cutoff shape and is not a group-invariant substitute.

There is a useful convention warning. Some authors impose $g(e_0)=e_0$. We
write this subgroup as
\begin{equation}
 G^1=W^*\rtimes\Sp_{2n}(\mathbb C),
\end{equation}
which is obtained by setting $a=1$. The full stabilizer $G$ and $G^1$ give
different answers below and should not be conflated.

\section{A structural identity and what it sees}

Block triangularity gives, for every scalar $t$,
\begin{equation}
 \det(I-tg)=(1-at)\det(I-tA),\qquad
 \Spec(g)=\Spec(A)\cup\{a\}.
\end{equation}
The shear coordinate $\ell$ can change Jordan structure but is invisible to
the characteristic polynomial. Therefore any probability measure $\nu$ for
which the coefficientwise moments exist satisfies
\begin{equation}
 \int_G\prod_i\det(I-t_i g)^{-1}\,d\nu(g)
 =\int_{\Sp_{2n}(\mathbb C)\times\mathbb C^\times}
 \prod_i\frac{\det(I-t_iA)^{-1}}{1-at_i}\,d\bar\nu(A,a),
\end{equation}
where $\bar\nu$ is the pushforward to the Levi quotient. This gives a useful
diagnostic for heat-kernel or user-chosen noncompact probabilities: the answer
depends on the chosen measure, but never on its conditional distribution in
the unipotent coordinate.

The code tests (4) modulo $q$ on deterministic samples from every constructed
finite group. It is important that this structural test is separate from the
permutation-representation test in Section~\ref{sec:finite}.

\section{Two compact replacements}

A maximal compact subgroup of the full stabilizer is
\[
 K=\Sp(n)\times U(1),\qquad V|_K=W\oplus\chi,
\]
where $\chi(z)=z$ on the radical line. Write
\[
 \Sfun_{H,X}(\mathbf t)=\int_H\prod_i\det(I-t_i h|_X)^{-1}\,dh
\]
for a compact group with normalized Haar probability. Direct-sum
multiplicativity and product Haar measure give
\[
 \Sfun_{K,V}=\Sfun_{\Sp(n),W}\Sfun_{U(1),\chi}.
\]
Since
\[
 \prod_i(1-zt_i)^{-1}=\sum_{d\ge0}z^d h_d(\mathbf t),
 \qquad \int_{U(1)}z^d\,dz=\delta_{d0},
\]
we obtain the exact identity
\begin{equation}
 \boxed{\Sfun_{K,V}=\Sfun_{\Sp(n),W}.}
\end{equation}
The radical line disappears because it has a nontrivial $U(1)$ weight.

For the fixed-radical group $G^1$, a maximal compact subgroup is $\Sp(n)$ and
$V|_{\Sp(n)}=W\oplus\mathbf 1$. The extra eigenvalue is now $1$, so
\begin{equation}
 \boxed{\Sfun_{\Sp(n),W\oplus\mathbf1}
 =\Sfun_{\Sp(n),W}\prod_i(1-t_i)^{-1}.}
\end{equation}

\subsection*{Exact rank-one check}

For $n=1$, restrict to the maximal torus of $SU(2)=\Sp(1)$ with weights
$z,z^{-1}$. Weyl integration is the exact constant-term operation
\[
 \int_{SU(2)}f=\operatorname{CT}_z\bigl((1-z^2)f(z)\bigr).
\]
The verification represents every complete homogeneous character as a Laurent
polynomial. It extracts both the $z$ and $U(1)$ constant terms and tests all 30
partitions through total degree six. Representative results are:
\begin{center}
\begin{tabular}{@{}lrrrr@{}}
\toprule
$\tau$ & $K$ direct & $\Sp(1),W$ & $G^1$ compact direct & formula (7)\\
\midrule
$()$      &1&1&1&1\\
$(1)$     &0&0&1&1\\
$(2)$     &0&0&1&1\\
$(1,1)$   &1&1&2&2\\
$(2,1)$   &0&0&2&2\\
$(2,2)$   &1&1&3&3\\
\bottomrule
\end{tabular}
\end{center}
All 30 comparisons pass exactly, with integer arithmetic.

\section{The full algebraic invariant series}

Even though Haar averaging fails, the covariant invariant series is meaningful:
\begin{equation}
 \Sfun^{\mathrm{inv}}_{G,V}
 =\sum_\tau\dim\left(\bigotimes_j\Sym^{\tau_j}V\right)^Gm_\tau.
\end{equation}

\textbf{Theorem.} For the full stabilizer $G$ in (1),
\begin{equation}
 \boxed{\Sfun^{\mathrm{inv}}_{G,V}=1.}
\end{equation}

\textit{Proof.}
Embed $\bigotimes_j\Sym^{\tau_j}V$ into $V^{\otimes d}$, where
$d=|\tau|$. The subgroup $\operatorname{diag}(I_W,a)$ acts with weight $a^r$
on tensors having $r$ radical-line factors. Its invariant subspace is thus
contained in $W^{\otimes d}$.

For $\ell\in W^*$, differentiate the action of
\[
 u_\ell=\begin{pmatrix}I_W&0\\ \ell&1\end{pmatrix}.
\]
On $x\in W^{\otimes d}$ the derivative is a sum of $d$ contractions, one for
each tensor position, and the output of the $j$th contraction has its unique
$e_0$ in position $j$. These outputs lie in distinct direct summands. Hence
invariance forces every positional contraction against every $\ell$ to vanish.
Already at the first position, the contractions by a basis of $W^*$ recover
all coordinates of $x$, so $x=0$. Thus there are no positive-degree
invariants. \hfill$\square$

This proof gives a compact computational certificate: after the torus weight
filter, stack the first-position contractions against a basis of $W^*$. The
result is a coordinate permutation of rank $(2n)^d$. The implementation records
full rank for $n=1,2$ and $1\le d\le4$ (eight checks), so every recorded kernel
has dimension zero. This certificate is stronger and cheaper than building
large Reynolds matrices for a nonreductive group.

\section{Finite odd symplectic groups and a genuine complex representation}
\label{sec:finite}

Let $E=W\oplus L$ over a prime field $\mathbb F_q$. The same block proof gives
\begin{equation}
 G^{\mathrm{odd}}_{2n+1}(q)
 \cong W^*\rtimes\bigl(\Sp_{2n}(q)\times\mathbb F_q^\times\bigr),
\end{equation}
and hence
\begin{align}
 |G^{\mathrm{odd}}_{2n+1}(q)|
 &=q^{2n}(q-1)|\Sp_{2n}(q)|,\notag\\
 |\Sp_{2n}(q)|&=q^{n^2}\prod_{j=1}^n(q^{2j}-1).
\end{align}

Matrices with entries in $\mathbb F_q$ are not automatically complex matrices.
We therefore use the honest complex permutation representation on
$S=\mathbb P(E)$. If $P_g$ is its permutation matrix and $c_r(g)$ is the
number of $r$-cycles, then
\begin{equation}
 \det(I-tP_g)^{-1}=\prod_{r\ge1}(1-t^r)^{-c_r(g)}.
\end{equation}
Consequently the proposed finite formula is
\begin{equation}
 \boxed{\Sfun_{G(q),\mathbb C[S]}(\mathbf t)
 =\frac1{|G(q)|}\sum_{g\in G(q)}
 \prod_i\prod_{r\ge1}(1-t_i^r)^{-c_r(g)}.}
\end{equation}
For a partition $\tau$, Burnside's lemma also gives
\begin{equation}
 [m_\tau]\Sfun_{G(q),\mathbb C[S]}
 =\#\left(G(q)\backslash\prod_j\operatorname{MSet}_{\tau_j}(S)\right).
\end{equation}

There are exactly two orbits on $\mathbb P(E)$: the radical point and all
nonradical points. The latter assertion follows because $\Sp(W)$ is transitive
on nonzero vectors of $W$, while $\ell(w)$ adjusts the radical coordinate.
Therefore
\begin{equation}
 \boxed{[m_{(1)}]\Sfun_{G(q),\mathbb C[\mathbb P(E)]}=2.}
\end{equation}

\section{Independent verification algorithms}

For each $(n,q)$, the program performs the following steps.
\begin{enumerate}[leftmargin=*,itemsep=3pt]
\item Enumerate every $2n\times2n$ matrix $A$ satisfying
$A^{\mathsf T}JA=J$, then every triple $(A,\ell,a)$ in (10). Check the observed
orders against (11), preservation of the degenerate form, and deterministic
closure samples.
\item Enumerate normalized representatives of every projective point and
construct the permutation $P_g$ for every group matrix.
\item \textbf{Direct matrix route:} compute $p_k=\operatorname{tr}(P_g^k)$
from explicit permutation powers and use Newton's recurrence
\[
 d h_d=\sum_{k=1}^d p_kh_{d-k}.
\]
Average $\prod_jh_{\tau_j}$ over all matrices.
\item \textbf{Cycle route:} independently find the cycles and expand the
right side of (12) by dynamic programming. Average the resulting products and
compare every coefficient with the direct route.
\item \textbf{Orbit route:} in the two smaller cases, enumerate the orbits on
multisets for $\tau=(1),(2),(1,1)$ and compare with (14).
\end{enumerate}

The direct route does not call the cycle-product routine. This separation
matters: otherwise (13) would only be tested against a rearrangement of itself.

\section{Exact results}

\begin{center}
\begin{tabularx}{\linewidth}{@{}rrrrrrX@{}}
\toprule
$n$ & $q$ & $\dim E$ & $|\Sp_{2n}(q)|$ & $|G(q)|$ & $|\mathbb P(E)|$ & tested degrees\\
\midrule
1&2&3&6&24&7&all partitions through 4\\
1&3&3&24&432&13&all partitions through 4\\
1&5&3&120&12,000&31&all partitions through 3\\
2&2&5&720&11,520&31&all partitions through 3\\
\bottomrule
\end{tabularx}
\end{center}
All four form, order, radical-point, determinant, closure, and coefficient
checks pass. In total, 38 matrix-average coefficients match 38 cycle-formula
coefficients exactly. Selected values illustrate both stability and genuine
dimension dependence:
\begin{center}
\begin{tabular}{@{}ccrrrrrr@{}}
\toprule
$(n,q)$ & $m_{(1)}$ & $m_{(2)}$ & $m_{(1,1)}$ & $m_{(3)}$ & $m_{(2,1)}$ & $m_{(1,1,1)}$\\
\midrule
$(1,2)$&2&5&6&11&18&25\\
$(1,3)$&2&5&6&12&19&28\\
$(1,5)$&2&5&6&14&22&30\\
$(2,2)$&2&6&7&19&31&43\\
\bottomrule
\end{tabular}
\end{center}
For $(n,q)=(1,2)$ and $(1,3)$, explicit orbit enumeration gives respectively
$2$, $5$, and $6$ orbits for $\tau=(1),(2),(1,1)$, exactly matching both
averaging routes.

\section{Interpretation, limitations, and new tools}

The results separate four notions that are easy to blur:
\begin{itemize}[leftmargin=*,itemsep=3pt]
\item The full complex-group Haar expression is undefined, and no computation
can validate it without first choosing a new measure.
\item Maximal-compact averaging is canonical up to conjugacy, but it forgets
the unipotent radical. For the full stabilizer it also removes the radical
line by $U(1)$ weight; for $G^1$ it retains that line as a trivial summand.
\item The full covariant invariant series is well-defined but collapses to $1$
in the defining representation. Interesting odd symplectic representation
theory therefore requires other modules, mixed $V,V^*$ constructions, or
polynomial-function invariants.
\item Finite-field permutation actions give exact, nontrivial, reproducible
$\Lambda$-distributions, but they are different representations of different
finite groups and should not be presented as complex matrix realizations of
the defining modular matrices.
\end{itemize}

Two reusable tools emerge. First, the \emph{Levi visibility test} (4)--(5)
immediately decides which nonreductive coordinates an eigenvalue statistic can
detect. Second, the \emph{torus-filter plus derivative certificate} proves
invariant vanishing using small full-rank contraction maps instead of a
nonexistent Reynolds average.

\section{Reproducibility}

From the repository root, run
\begin{verbatim}
.venv/bin/python \
  lambda_distributions/proofs/proctor_odd_symplectic_group/verification.py
\end{verbatim}
and open the interactive report with
\begin{verbatim}
.venv/bin/marimo edit \
  marimo_notebooks/proctor_odd_symplectic.py
\end{verbatim}
The verification uses only exact integer and rational arithmetic. Exhaustive
group enumeration is intentionally limited to prime fields and to the four
displayed cases; extending to larger ranks should replace all-matrix search by
standard symplectic generators.

\begin{thebibliography}{9}
\bibitem{Okada}
S.~Okada, \emph{A bialternant formula for odd symplectic characters and its
application}, 2019, \url{https://arxiv.org/abs/1905.12964}.
\bibitem{Burnside}
W.~Burnside, \emph{Theory of Groups of Finite Order}, second edition, 1911.
\bibitem{Haar}
A.~Haar, \emph{Der Massbegriff in der Theorie der kontinuierlichen Gruppen},
Annals of Mathematics 34 (1933), 147--169.
\end{thebibliography}

\end{document}
